\section{Coroutine dan Event}

\begin{learningoutcomes}
\item Memahami konsep coroutine dan IEnumerator
\item Menggunakan yield return untuk delay
\item Mengimplementasikan event system
\end{learningoutcomes}

\subsection{Coroutine}

\begin{definisi}[Coroutine]
Coroutine adalah fungsi yang dapat menjeda eksekusinya sendiri dan melanjutkan di frame berikutnya \cite{UnityScriptingAPI2024}.
\end{definisi}

\begin{lstlisting}[language={[Sharp]C}]
using System.Collections;
using UnityEngine;

IEnumerator Countdown()
{
    for (int i = 5; i > 0; i--)
    {
        Debug.Log(i);
        yield return new WaitForSeconds(1f);
    }
}
\end{lstlisting}

\subsection{Event dan Delegate}

\begin{lstlisting}[language={[Sharp]C}]
using System;

public static event Action OnGameOver;
public static event Action<int> OnScoreChanged;

// Subscribe di OnEnable, Unsubscribe di OnDisable
void OnEnable() { GameManager.OnGameOver += HandleGameOver; }
void OnDisable() { GameManager.OnGameOver -= HandleGameOver; }
\end{lstlisting}

\begin{catatan}
Selalu unsubscribe event di OnDisable untuk menghindari memory leaks.
\end{catatan}
